\section{Introduction to Purely Modal Languages}\label{pure modal languages section}

\par Purely modal languages are modal languages that remove the use of the disjunction (and consequently the conjunction) and $\top$ (and consequently $\bot$). The pure modal language we investigate in this chapter is able to express everything in the basic modal language. We prove this in Proposition~\ref{pure modal expresses basic proposition}. 
%Ignored for First submission
%\nnote{(The above in this paragraph has been redone. Please check.)} 
For this, we define the similarity type that will be used (Definition~\ref{similarity type}) to be arbitrary, including the terms $\iota_0, \iota_1,$ and $\iota_2$, which are nullary, unary, and binary terms, respectively. For the language we are considering, we also allow modal terms to be composed under certain restrictions. This language is taken from~\cite{goranko2006elementary}. Also, the definitions in this section are adapted from~\cite{goranko2006elementary} as well. For our study, we restrict the language such that $\iota_0$ is the only nullary modality we will consider. In this chapter, when we say $\tau$ is a similarity type, we mean that $\tau$ is a similarity type that includes the terms $\iota_0, \iota_1,$ and $\iota_2$.



\begin{definition}

    \textbf{(The set $\mtt$)} Let $\tau$ be a similarity type. We define the set $\mtt$ (model terms from $\tau$) inductively as follows:

    \begin{enumerate}[label=(\roman*)]

        \item if $\alpha \in \tau$ then $\alpha\in\mtt$,

        \item if $m > 0 \text{ and } \alpha, \beta_1, \ldots, \beta_m \in \mtt \text{ and } \rho(\alpha) = m$ then $\alpha(\beta_1,\ldots,\beta_m) \in \mtt \text{ and } \rho(\alpha(\beta_1,\ldots$\\
        $,\beta_m)) = \rho(\beta_1) + \cdots + \rho(\beta_m)$,

        \item and nothing else.

    \end{enumerate}

    %\wnote{This definition deviates from Gonranko and Vakarelov in that it does not include constant formulas as 0-ary terms. Why? If there is a reason for this deviation, it should be discussed and motivated. }

\end{definition}

\par Note: It is suitable to mention here the inductive definition of $\iota_n$. We have the modal terms $\iota_0$, $\iota_1$, $\iota_2$ in our similarity type, and so in $\mtt$, we define $\iota_n = \iota_2(\iota_1, \iota_{n-1})$ $(n\geq3)$. We call the modal terms satisfying (i) atomic modal terms and those that satisfy (ii) composite modal terms.
\par Next, we define the frames we will be using. Note that if for any $(n+1)$-ary relations we will denote them as $R_\alpha (wu_1\ldots u_n)$ or $R_\alpha wu_1\ldots u_n$.
%Ignored for First submission
%\nnote{Phrasing?}

\begin{definition}\label{mtt-frame}

    \textbf{($\mtt$-frame)} Let $\tau$ be a model similarity type, then a \textit{$\mtt$-frame} is a tuple $\Fmtt = (W, \{ R_{\alpha} \}_{\alpha \in \mtt})$ that consists of:

    \begin{enumerate}[label=(\roman*)]

        \item A non-empty set $W$ (of worlds).

        \item For each $m\geq0$, each $m$-ary modal term $\alpha$ in $\mtt$, an $(m+1)$-ary relation $R_{\alpha}$, defined as follows:

        \begin{itemize}

            \item $R_{\iota_0} = W \coms R_{\iota_1} = \{ (w,w) \mid w \in W \} \coms R_{\iota_2} = \{ (w,w,w) \mid w\in W\}$,

            \item for every modal term $\alpha \in\tau$, we have $R_\alpha \subseteq W^{\rho(\alpha)+1}$,

            \item letting $\rho(\beta_i)=b_i, 1\leq i \leq m$, $R_{\alpha(\beta_1,\ldots,\beta_m)} = \{ (w, w^1_1, \ldots, w^1_{b_1}, w^2_{1},\ldots, w^2_{b_2},\ldots, w^m_{1},\ldots, w^m_{b_m}) \in W^{b_1 + b_2 + \cdots + b_m + 1} \mid \text{there } \text{are } u_1, \ldots, u_m \in W$ where $R_{\alpha}wu_1\ldots u_m\;\text{and}\;R_{\beta_i}u_iw^i_1\ldots w^i_{b_i}$ for all $1\leq i\leq m$\} = $R_\alpha\circ (R_{\beta_1},\ldots,R_{\beta_m})$, 

            \item and nothing else.

        \end{itemize}

    \end{enumerate}
    Note that when it is clear from the context, we often omit the subscript $\mtt$ from $\F_{\mtt}$.
\end{definition}

\par

\begin{definition}

    \textbf{($\mtt$-model)} Let $\Fmtt$ be a $\mtt$-frame. A \textit{$\mtt$-model} (denoted $\Mmtt$) is an ordered pair $(\Fmtt,V)$, where $V:\PVAR\to\mathcal{P}(W)$ is a valuation on the set of propositional variables \PVAR. We denote a $\mtt$-model as $\Mmtt = (W,\{ R_{\alpha} \}_{\alpha \in \mtt},V)$, but again, we will often omit the subscript $\mtt$ when it's clear from the context.

\end{definition}



\par Now we introduce the purely modal language.



\begin{definition}

    \textbf{(Purely modal language)}  Let $\tau$ be a similarity type. The formulas in the purely modal language (which we denote as $\Lp$) are given by the rule:

    \begin{equation*}
        \varphi := p\mid \neg\varphi \mid [\alpha](\varphi_1,\ldots,\varphi_m)
    \end{equation*}

    \par where $p\in$ \PVAR, $\alpha \in \mtt$ and $\rho(\alpha)=m$.

\end{definition}

Thus, the purely modal language is constructed from a set $\PVAR$ of propositional variables, negated formulas, and what we call boxed formulas.



\begin{definition}\label{dual of alpha}

    \textbf{(Dual of $\alpha$)} For any modal formula $[\alpha](\phi_1,\phi_2,\ldots,\phi_{\rho(\alpha)})$, we define its \textit{dual} as $\langle\alpha\rangle(\phi_1,\phi_2,\ldots,\phi_{\rho(\alpha)}) := \neg[\alpha](\neg\phi_1,\neg\phi_2,\ldots,\neg\phi_{\rho(\alpha)})$.

\end{definition}



\begin{definition}\label{satisfiability clauses - pure}

    \textbf{(Semantics)} Let $\Mmtt = (W,\{ R_{\alpha} \}_{\alpha \in \mtt},V)$ be a $\mtt$-model and $w\in W$, then

    \begin{itemize}

        \item $\Mmtt , w \Vdash p$ \quad iff \quad $p\in V(p)$

        \item $\Mmtt , w \Vdash \neg\phi$ \quad iff \quad $\Mmtt , w \not\Vdash \phi$

    \end{itemize}    

    \par Now for the boxed formulas, we have two cases: $\rho(\alpha)=m > 0$ and $\rho(\alpha) = 0$.

    \par So, if $\rho(\alpha)=m > 0$

    \begin{center}

        $\Mmtt,  w \Vdash [\alpha](\phi_1,\phi_2,\ldots,\phi_m)$\\

        iff\\

        for all $v_1,v_2, \ldots, v_m \in W$ with $R_{\alpha}wv_1v_2\ldots v_m$ we have\\

        $\Mmtt, v_1\Vdash\phi_1$ or $\Mmtt, v_2\Vdash\phi_2$ or $\cdots$ or $\Mmtt, v_m\Vdash\phi_m$

    \end{center}







\par Now for $\iota_0$, (i.e., when $m=0$)

    \begin{center}

        $\Mmtt,  w \Vdash [\iota_0]$ \quad iff \quad $w\in W$

    \end{center}





    %prev definition of when m=0

% \par Now for $m=0$ then

%     \begin{center}

%         $\Mmtt,  w \Vdash [\alpha]$ \quad iff \quad $w\in R_\alpha$

%     \end{center}

% \par and its dual as:

% \begin{center}

%         $\Mmtt,  w \Vdash \langle\alpha\rangle$ \quad iff \quad $w\not\in R_\alpha$

%     \end{center}

\end{definition}

\begin{remark}

\par Now for the dual of the box ($m > 0$) using Definition~\ref{dual of alpha}:

    \begin{center}

        $\Mmtt,  w \Vdash \langle\alpha\rangle(\phi_1,\phi_2,\ldots,\phi_m)$\\

        iff\\

        there are states $v_1,v_2, \ldots, v_m \in W$ with $R_{\alpha}wv_1v_2\ldots v_m$ where \\

        $\Mmtt, v_i\Vdash\phi_i$ for each $i \in \{1,2,\ldots,m \}$

    \end{center}

\par and for $m=0$:

    \begin{center}

        $\Mmtt,  w \Vdash \langle\iota_0\rangle$ \quad iff \quad $w\not\in W$

    \end{center}
\end{remark}



\par Note that from these definitions it follows that

\begin{center}

    $\lsb\alpha\rsb(\lsb\beta_1\rsb(\phi_1^1\text{,\,}\phi_2^1\text{,\,}\ldots \text{,\,}\phi_{\rho(\beta_1)}^1)\text{,\,} \ldots \text{,\,} \lsb\beta_{m}\rsb(\phi_1^m \text{,\,}\ldots \text{,\,}\phi_{\rho(\beta_m)}^m))$

\end{center}

\par is equivalent to 

\begin{center}

    $\lsb\alpha(\beta_1\text{,\,}\ldots\text{,\,}\beta_{\rho(\alpha)})\rsb(\phi_1^1\text{,\,}\phi_2^1\text{,\,}\ldots \text{,\,}\phi_{\rho(\beta_1)}^1\text{,\,} \phi_1^2\text{,\,} \ldots \text{,\,}\phi_{\rho(\beta_2)}^2\text{,\,} \ldots\text{,\,} \phi_1^m \text{,\,}\ldots \text{,\,}\phi_{\rho(\beta_m)}^m)$.

\end{center}



\par Using the above equivalence, we can define a subset of formulas from the pure modal language, which we will call ``Maximally composed" formulas. Where we compose as many boxes (or diamonds) as possible. We call this subset $\LpMax$.



\begin{definition}\label{maximally composed formulas definition - pure}
    \textbf{(Maximally composed formulas)} The set of maximally composed $\Lp$ formulas ($\LpMax$) is defined by the union of the sets of formulas of the forms $\phi$ and $\psi$ given by the mutually inductive rule:

    \begin{center}

        $\phi := p \mid \neg\psi \mid [\alpha](\psi_1, \psi_2, \ldots,\psi_{\rho(\alpha)})$\\

        $\psi := p \mid \neg\phi \mid \langle\alpha\rangle(\phi_1, \phi_2, \ldots,\phi_{\rho(\alpha)})$

    \end{center}\

\par where $p\in\PVAR$ and $\alpha\in\mtt$.

\end{definition}



\par The rest of this section describes how we can capture the semantics of the disjunction in $\Lp$.



\paragraph{Capturing the semantics of the disjunction:}

\par Every similarity type $\tau$ contains the modal terms $\iota_0, \iota_1,$ and $\iota_2$ and so all $\mtt$-models contain the respective corresponding relations. Using the definition that $\iota_n = \iota_2(\iota_1, \iota_{n-1})$ $(n\geq3)$, we can then conclude that the $\riotn$ relation is $R_{\iota_n} = \{ (\underbrace{w,\ldots,w}_{n+1 \text{times}}) \mid w\in W\}$. So, for every $\mtt$-model $\M$:



% \begin{definition}

%     \textbf{The iota-n ($\iota$) relation}: For any $\tau$-model $\M$ we define:

% $R_{\iota_0}$ to be: $R_{\iota_0} := W$, $R_{\iota_1}$ as: $R_{\iota_1} := \{ (w,w) \mid w\in W\}$,

% $R_{\iota_2}$ as: $R_{\iota_2} := \{ (w,w,w) \mid w\in W\}$, Now using the definition of $\iota_n := \iota_2(\iota_1,\iota_{n-1})$ we have that $R_{\iota_n} = \{ (\underbrace{w,\ldots,w}_{n+1 \text{times}}) \mid w\in W\}$. Note that $\rho(\iota_n) = n$.

% \end{definition}



\begin{center}

    $\M, w \Vdash [\alpha](\phi_1, \phi_2, \ldots, \phi_{\rho(\alpha)})$\\

    iff\\

    for all $v_1,v_2, \ldots, v_{\rho(\alpha)}$ with $R_{\alpha}wv_1v_2\ldots v_{\rho(\alpha)}$ it is the case that\\

    $\M, v_1\Vdash\phi_1$ or $\M, v_2\Vdash\phi_2$ or $\cdots$ or $\M, v_{\rho(\alpha)}\Vdash\phi_{\rho(\alpha)}$.

\end{center}



\par This is for an arbitrary relation $R_\alpha$, but if we use the $\iota$-relation we get:

\begin{center}

    $\M, w \Vdash [\iota_n](\phi_1, \phi_2, \ldots, \phi_{n})$\\

    iff\\

    for all $v_1,v_2, \ldots, v_{n}$ with $R_{\iota_n}wv_1v_2\ldots v_{n}$ it is the case that\\

    $\M, v_1\Vdash\phi_1$ or $\M, v_2\Vdash\phi_2$ or $\cdots$ or $\M, v_{n}\Vdash\phi_{n}$.

\end{center}

\par Now by definition, if $R_{\iota_n}wv_1v_2\ldots v_{n}$, then $(w,v_1,\ldots,v_n) \in R_{\iota_n}.$ This means that  $w=v_1=\cdots=v_n$. So then:

\begin{center}

    $\M, w \Vdash [\iota_n](\phi_1, \phi_2, \ldots, \phi_{n})$\\

    iff\\

    %for all $y_1,y_2, \ldots, y_{n}$ with $R_{\iota_n}xy_1y_2\ldots y_{n}$ we have\\

    $\M, w\Vdash\phi_1$ or $\M, w\Vdash\phi_2$ or $\cdots$ or $\M, w\Vdash\phi_{n}$\\

    iff\\

    $\M, w \Vdash \phi_1 \vee \phi_2 \vee \cdots \vee \phi_n$.

\end{center}



\par As we can see, using the correct modality, we can capture the semantics of a disjunction. Now, note for the $\iota_1$ and $\iota_0$ we have that:

\begin{center}

    $\M, w \Vdash [\iota_1](\phi_1)$\\

    iff\\

    $\M, w \Vdash \phi_1$

\end{center}

\par and we have

\begin{align*}
    \M, w \Vdash [\iota_0] \quad &\text{iff} \quad x\in R_{\iota_0} = W\\
    &\text{iff} \quad \text{always}\\
    &\text{iff} \quad \M, w \Vdash \top.
\end{align*}

\par It may seem that $\iota_1$ doesn't play much of a role, but it is useful since any formula $\phi$ is equivalent to $[\iota_1]\phi$. This equivalence will be used later. We also note that $\iota_0$ captures the semantics of $\top$, which is used in the basic modal language.

\par Next, we want to prove that every modal language $ML_\tau$ is expressively equivalent to a purely modal language $\Lp$. To do this, we will need to use induction on $\Lp$, but due to how we can compose terms in $\mtt$, it is helpful to introduce functions that measure the complexity of a formula. We construct such a function next.

\begin{definition}\label{cep of modal term def}
    \textbf{(Cumulative Depth of a modal term)} Let $\tau$ be a similarity type. We define the function $cep:\mtt\to\mathbb{N}$ as follows:
    \begin{enumerate}
        \item if $\alpha\in\mtt$ and $\alpha$ is atomic then $cep(\alpha)=1$.
        \item  and if $\alpha\in\mtt$ and $\alpha$ is composite, i.e., $\alpha=\pi(\beta_1,\ldots,\beta_{\rho(\pi)})\in\mtt$ (for some $\pi,\beta_1,\ldots,\beta_{\rho(\pi)}\in\mtt$), then $cep(\alpha)=1 + cep(\pi) + \Sigma_{i=1}^{\rho(\pi)}cep(\beta_i)$.
    \end{enumerate}
\end{definition}

\begin{definition}\label{com of formula def}
    \textbf{(Complexity of a formula)} We let $\tau$ be a similarity type. We define the function $com: \Lp\to \mathbb{N}$ as follows ($m=\rho(\alpha)$):
    \begin{enumerate}
        \item if $\phi = p$ then $com(\phi)=0$
        \item if $\phi = \neg\psi$ then $com(\phi) = 1 + com(\psi)$
        \item if $\phi = [\alpha](\psi_1,\ldots,\psi_{m})$ then $com(\phi) = cep(\alpha) + \Sigma_{i=1}^{m} com(\psi_i)$.
    \end{enumerate}
\end{definition}

\begin{lemma}\label{composed is higher com lemma}
    Let $com:\Lp\to\mathbb{N}$ be as defined above and let $\pi(\beta_1,\ldots,\beta_m)\in\mtt$ with $\rho(\pi) = m$ and $\rho(\beta_i) = b_i$ for all $1\leq i\leq m$. Then $com([\pi(\beta_1,\ldots,\beta_m)](\phi_1^1,\ldots,\phi_{b_1}^1,\ldots,\phi_1^m,\ldots,\phi_{b_m}^m))
    \\ > com([\pi]([\beta_1](\phi_1^1,\ldots,\phi_{b_1}^1),\ldots,[\beta_m](\phi_1^m,\ldots,\phi_{b_m}^m)))$.
\end{lemma}
    
\begin{proof}
    Let $com:\Lp\to\mathbb{N}$ be as defined above and let $\pi(\beta_1,\ldots,\beta_m)\in\mtt$ with $\rho(\pi) = m$ and $\rho(\beta_i) = b_i$ for all $1\leq i\leq m$. Now,
    \begin{align*}
        com([\pi(\beta_1,\ldots,\beta_m)](\phi_1^1,\ldots,\phi_{b_1}^1,&\ldots,\phi_1^m,\ldots,\phi_{b_m}^m))\\
        &= cep(\pi(\beta_1,\ldots,\beta_m)) + \Sigma_{i=1}^m\Sigma_{j=1}^{b_i} com(\phi_j^i)\\
        & = 1 + cep(\pi) + \Sigma_{i=1}^{m}cep(\beta_i) + \Sigma_{i=1}^m\Sigma_{j=1}^{b_i} com(\phi_j^i),
    \end{align*}
    
    and
    \begin{align*}
        com([\pi]([\beta_1](\phi_1^1,\ldots,\phi_{b_1}^1),\ldots &,[\beta_m](\phi_1^m,\ldots,\phi_{b_m}^m)))\\
        & = cep(\pi) + \Sigma_{i=1}^{m}com[\beta_i](\phi_1^i,\ldots,\phi_{b_i}^i)\\
        & = cep(\pi) + \Sigma_{i=1}^{m}(cep(\beta_i) + \Sigma_{j=1}^{b_i} com(\phi_j^i))\\
        & = cep(\pi) + \Sigma_{i=1}^{m}cep(\beta_i) + \Sigma_{i=1}^m\Sigma_{j=1}^{b_i} com(\phi_j^i).
    \end{align*}
    Hence, our original statement is satisfied.
\end{proof}

\begin{lemma}\label{composite written with leading atomic lemma}
    Every modality $\alpha\in\mtt$ can be rewritten such that the leading modality is atomic. 
    %Ignored for First submission
%\nnote{Leading modality is not defined. Is it common knowledge, or should I define it? What would a definition look like?}
\end{lemma}

\begin{proof}
    We do this by induction on $\mtt$. \textit{Base case:} If $\alpha$ is atomic, we are done. \textit{Induction hypothesis:} Assume that the statement is true for any modality $\alpha$ such that $cep(\alpha)<m$. \textit{Induction step:} Let $\alpha$ be composite such that $cep(\alpha) = m$. Now $\alpha = \pi(\beta_1,\ldots,\beta_{\rho(\pi)})$
    for some $\pi,\beta_1,\ldots,\beta_{\rho(\pi)}\in\mtt$. If $\pi$ is atomic, we are done. If $\pi$ is composite, since $cep(\pi)<cep(\alpha)$, we can rewrite $\pi$ as $\pi = \lambda(\mu_1,\ldots,\mu_{\rho(\lambda)})$ where $\lambda$ is atomic using our induction hypothesis.
    Now $\alpha = (\lambda(\mu_1,\ldots,\mu_{\rho(\lambda)}))(\beta_1,\ldots,\beta_{\rho(\pi)})$, but we can reindex the $\beta_i$'s (letting $\rho(\lambda) = l$ and $\rho(\mu_i) = m_i$ for all $1\leq i\leq l$) to get $\alpha = (\lambda(\mu_1,\ldots,\mu_l))(\beta_1^1,\ldots,\beta_{m_1}^1,\ldots,\beta_1^l,\ldots,\beta_{m_l}^l) = \lambda(\mu_1(\beta_1^1,\ldots,\beta_{m_1}^1),\ldots,\mu_l(\beta_1^l,\ldots,\beta_{m_l}^l))$. Hence, by induction on $\mtt$, every modality $\alpha\in\mtt$ can be rewritten such that the leading modality is atomic.
\end{proof}

\begin{proposition}\label{pure modal expresses basic proposition}
    Every modal language $ML_\tau$ is expressively equivalent to a purely modal language $\Lp$. 
\end{proposition}

\begin{proof}
    We begin by defining two translations $\zeta: ML_\tau \to\Lp$ and $\eta:  \Lp \to ML_\tau$ by induction.
    \begin{enumerate}
        \item If $\phi = p\in\PVAR$ then $\zeta(\phi) = p$
        \item If $\phi = \top$ then $\zeta(\phi) = [\iota_0]$
        \item If $\phi = \neg\psi$ then $\zeta(\phi) = \neg\zeta(\psi)$
        \item If $\phi = \psi_1\vee\psi_2$ then $\zeta(\phi) = [\iot{2}](\zeta(\psi_1),\zeta(\psi_2))$
        \item If $\phi = [\alpha](\psi_1,\ldots,\psi_{\rho(\alpha)})$ then $\zeta(\phi) = [\alpha](\zeta(\psi_1),\ldots,\zeta(\psi_{\rho(\alpha)}))$
    \end{enumerate}
    and
    \begin{enumerate}
        \item If $\phi = p\in\PVAR$ then $\eta(\phi) = p$
        \item If $\phi = \neg\psi$ then $\eta(\phi) = \neg\eta(\psi)$
        \item If $\phi = [\iot{0}]$ then $\eta(\phi) = \top$
        \item If $\phi = [\iot{1}]\psi$ then $\eta(\phi) = \eta(\psi)$
        \item If $\phi = [\iot{2}](\psi_1,\psi_2)$ then $\eta(\phi) = \eta(\psi_1) \vee \eta(\psi_2)$
        \item If $\phi = [\alpha](\psi_1,\ldots,\psi_{\rho(\alpha)})$ and $\alpha$ is atomic then $\eta(\phi) = [\alpha](\eta(\psi_1),\ldots,\eta(\psi_{\rho(\alpha)}))$
        \item If $\phi = [\alpha](\psi_1,\ldots,\psi_{\rho(\alpha)})$ and $\alpha$ is composite then write $\alpha$ as $\pi(\beta_1,\ldots,\beta_{\rho(\pi)})$ where $\pi$ is atomic (Lemma~\ref{composite  written with leading atomic lemma}). After reindexing (letting $\rho(\pi) = m$ and $\rho(\beta_i) = b_i$) $\phi=[\pi(\beta_1,\ldots,\beta_m)](\psi_1^1,\ldots,\psi_{b_1}^1,\ldots,\psi_1^m,\ldots,\psi_{b_m}^m)$ then $\eta(\phi) = [\pi](\eta([\beta_1](\psi_1^1,\ldots,\psi_{b_1}^1)),\ldots,
        \\\eta([\beta_m](\psi_1^m,\ldots,\psi_{b_m}^m)))$.
    \end{enumerate}
    \par Note that every $\mtt$-model contains an $ML_\tau$-model as a reduct, obtained by omitting the relations obtained by composition and only retaining those corresponding to terms in $\tau\setminus\{\iot{0},\iot{1},\iot{2}\}$. We can therefore apply the definition of the semantics of $ML_\tau$ also to $\mtt$ models. Now we are left to prove two statements:
    \begin{enumerate}
        \item That for every $\mtt$-model $\M$ and state $w$ in $\M$, and for every formula $\phi\in ML_\tau$, it is the case that $\M,w\Vdash \phi$ iff $\M,w\Vdash \zeta(\phi)$,
        \item  and that for every $\mtt$-model $\M$ and state $w$ in $\M$, and for every formula $\phi\in\Lp$, it is the case that $\M,w\Vdash \phi$ iff $\M,w\Vdash \eta(\phi)$.
    \end{enumerate}
    \par We prove both of these by induction on $ML_\tau$ and $\Lp$ respectively. So, for the first statement, let $\M = (W,\{ R_{\alpha} \}_{\alpha \in \mtt},V)$ be a $\mtt$-model and $w$ a state, and let $\phi\in ML_\tau$.
    \par \textit{Case 1:} If $\phi = p \in\PVAR$ then $\M,w\Vdash p$ iff $\M,w\Vdash \zeta(\phi)$.
    \par \textit{Case 2:} If $\phi = \top$ then $\M,w\Vdash \top$ iff $w\in W$ iff $\M,w\Vdash [\iot{0}]$ iff $\M,w\Vdash \zeta(\phi)$.
    \par \textit{Case 3:} Assume the statement is true for $\psi$. If $\phi = \neg\psi$ then $\M,w\Vdash\phi$ iff $\M,w\Vdash \neg\psi$ iff $\M,w\not\Vdash \psi$ iff $\M,w\not\Vdash \zeta(\psi)$ iff $\M,w\Vdash \neg\zeta(\psi)$ iff $\M,w\Vdash \zeta(\phi)$.
    \par \textit{Case 4:} Assume the statement is true for $\psi_1$ and $\psi_2$. If $\phi = \psi_1\vee\psi_2$ then $\M,w\Vdash\phi$ iff $\M,w\Vdash\psi_1\vee\psi_2$ iff $\M,w\Vdash\psi_1$ or $\M,w\Vdash\psi_2$ iff $\M,w\Vdash\zeta(\psi_1)$ or $\M,w\Vdash\zeta(\psi_2)$ iff $\M,w\Vdash[\iot{2}](\zeta(\psi_1),\zeta(\psi_2))$ iff $\M,w\Vdash\zeta(\phi)$.
    \par \textit{Case 5:} Let $\alpha\in\tau$ and assume the statement is true for $\psi_1,\ldots,\psi_{\rho(\alpha)}$. If $\phi = [\alpha](\psi_1,\ldots,\psi_{\rho(\alpha)})$ then
    $\M,w\Vdash\phi$
    iff $\M,w\Vdash[\alpha](\psi_1,\ldots,\psi_{\rho(\alpha)})$
    iff for all states $v_1,\ldots,v_{\rho(\alpha)}$ such that $R_\alpha wv_1\ldots v_{\rho(\alpha)}$
    and we have that for some $1\leq i\leq \rho(\alpha)$, $\M,v_i\Vdash\psi_i$. This is true 
    iff $\M,v_i\Vdash\zeta(\psi_i)$ for some $1\leq i\leq \rho(\alpha)$
    iff $\M,w\Vdash[\alpha](\zeta(\psi_1),\ldots,\zeta(\psi_{\rho(\alpha)}))$ iff $\M,w\Vdash\zeta(\phi)$.
    
    \par For the second statement, let $\M = (W,\{ R_{\alpha} \}_{\alpha \in \mtt},V)$ be a $\mtt$-model and $w$ a state, and let $\phi\in \Lp$.
    \par \textit{Case 1:} If $\phi = p \in\PVAR$ then $\M,w\Vdash p$ iff $\M,w\Vdash \eta(\phi)$.
    \par \textit{Case 2:} Assume the statement is true for $\psi$. If $\phi = \neg\psi$ then $\M,w\Vdash\phi$ iff $\M,w\Vdash \neg\psi$ iff $\M,w\not\Vdash \psi$ iff $\M,w\not\Vdash \eta(\psi)$ iff $\M,w\Vdash \neg\eta(\psi)$ iff $\M,w\Vdash \eta(\phi)$.
    \par \textit{Case 3:} If $\phi = [\iot{0}]$ then $\M,w\Vdash\phi$ iff $\M,w\Vdash[\iot{0}]$ iff $w\in W$ iff $\M,w\Vdash \top$ iff $\M,w\Vdash \eta(\phi)$.
    \par \textit{Case 4:} Assume the statement is true for $\psi$. If $\phi = [\iot{1}]\psi$ then $\M,w\Vdash\phi$ iff $\M,w\Vdash[\iot{1}]\psi$ iff $\M,w\Vdash\psi$ (since $R_\iot{1}ww$ and $w$ is the only $\iot{1}$-successor) iff $\M,w\Vdash\eta(\psi)$ iff $\M,w\Vdash\eta(\phi)$.
    \par \textit{Case 5:} Assume the statement is true for $\psi_1$ and $\psi_2$. If $\phi = [\iot{2}](\psi_1,\psi_2)$ then $\M,w\Vdash\phi$ iff $\M,w\Vdash[\iot{2}](\psi_1,\psi_2)$ iff $\M,w\Vdash\psi_1$ or $\M,w\Vdash\psi_2$ iff $\M,w\Vdash\eta(\psi_1)$ or $\M,w\Vdash\eta(\psi_2)$ iff $\M,w\Vdash\eta(\psi_1) \vee \eta(\psi_2)$ iff $\M,w\Vdash\eta(\phi)$.
    \par \textit{Case 6:} Let $\alpha\in\mtt$ be atomic and assume the statement is true for $\psi_1,\ldots,\psi_{\rho(\alpha)}$. If $\phi = [\alpha](\psi_1,\ldots,\psi_{\rho(\alpha)})$ then $\M,w\Vdash\phi$ iff $\M,w\Vdash[\alpha](\psi_1,\ldots,\psi_{\rho(\alpha)})$ iff for all states $v_1,\ldots,v_{\rho(\alpha)}$ such that $R_\alpha wv_1\ldots v_{\rho(\alpha)}$ and $\M,v_i\Vdash\psi_i$, for some $1\leq i\leq \rho(\alpha)$. This is true iff $\M,v_i\Vdash\eta(\psi_i)$ for some $1\leq i\leq \rho(\alpha)$, iff $\M,w\Vdash[\alpha](\eta(\psi_1),\ldots,\eta(\psi_{\rho(\alpha)}))$ iff $\M,w\Vdash\eta(\phi)$.

    % \par \tcr{V1}
    % \par \textit{Case 7:} Let $\alpha\in\mtt$ be composite. If $\phi = [\alpha](\psi_1,\ldots,\psi_{\rho(\alpha)})$ then $\M,w\Vdash\phi$ iff $\M,w\Vdash[\alpha](\psi_1,\ldots,\psi_{\rho(\alpha)})$. Now, since $\alpha$ is composite, we can rewrite it as a modal term $\pi(\beta_1,\ldots,\beta_{\rho(\pi)})$ where $\pi$ is atomic (Lemma~\ref{composite written with leading atomic lemma}). After reindexing the $\psi$'s (letting $\rho(\pi) = m$ and $\rho(\beta_i) = b_i$) we get that $\M,w\Vdash[\alpha](\psi_1,\ldots,\psi_{\rho(\alpha)})$ is true iff $\M,w\Vdash[\pi(\beta_1,\ldots,\beta_m)](\psi_1^1,\ldots,\psi_{b_1}^1,\ldots,\psi_1^m,\ldots,\psi_{b_m}^m)$ iff $\M,w\Vdash[\pi]([\beta_1](\psi_1^1,\ldots,\psi_{b_1}^1),\ldots,[\beta_m](\psi_1^m,\ldots,\psi_{b_m}^m))$. This shows that $\phi = [\alpha](\psi_1,\ldots,\psi_{\rho(\alpha)})$ can be rewriten as $[\pi]([\beta_1](\psi_1^1,\ldots,\psi_{b_1}^1),\ldots,[\beta_m](\psi_1^m,\ldots,\psi_{b_m}^m))$ where $\pi$ is atomic. This then falls under Case 6 and so $\M,w\Vdash[\alpha](\psi_1,\ldots,\psi_{\rho(\alpha)})$ iff $\M,w\Vdash[\pi]([\beta_1](\psi_1^1,\ldots,\psi_{b_1}^1),\ldots,[\beta_m](\psi_1^m,\ldots,\psi_{b_m}^m))$
    
    \par \textit{Case 7:} Let $\alpha\in\mtt$ be composite. Let $\phi = [\alpha](\psi_1,\ldots,\psi_{\rho(\alpha)})$ and assume the statement is true for lower $com$ value (Definition~\ref{com of formula def}). Then $\M,w\Vdash\phi$ iff $\M,w\Vdash[\alpha](\psi_1,\ldots,\psi_{\rho(\alpha)})$. Now, since $\alpha$ is composite, we can rewrite it as a modal term $\pi(\beta_1,\ldots,\beta_{\rho(\pi)})$ where $\pi$ is atomic (Lemma~\ref{composite written with leading atomic lemma}). After reindexing the $\psi$'s (letting $\rho(\pi) = m$ and $\rho(\beta_i) = b_i$) we get that $\M,w\Vdash[\alpha](\psi_1,\ldots,\psi_{\rho(\alpha)})$ is true iff $\M,w\Vdash[\pi(\beta_1,\ldots,\beta_m)](\psi_1^1,\ldots,\psi_{b_1}^1,\ldots,\psi_1^m,\ldots,\psi_{b_m}^m)$ iff
    $$\M,w\Vdash[\pi]([\beta_1](\psi_1^1,\ldots,\psi_{b_1}^1),\ldots,[\beta_m](\psi_1^m,\ldots,\psi_{b_m}^m))$$ iff for any states $v_1,\ldots,v_m$ such that $R_\pi wv_1\ldots v_m$ we have $\M,v_i\Vdash[\beta_i](\psi_1^i,\ldots,\psi_{b_i}^i)$ for some $1\leq i \leq m$. Since $com([\beta_i](\psi_1^i,\ldots,\psi_{b_i}^i))<com(\phi)$, then by our induction hypothesis, this is true iff  $\M,v_i\Vdash\eta([\beta_i](\psi_1^i,\ldots,\psi_{b_i}^i))$ for some $1\leq i \leq m$, iff $$\M,w\Vdash[\pi](\eta([\beta_1](\psi_1^1,\ldots, \psi_{b_1}^1)),\ldots,\eta([\beta_m](\psi_1^m, \ldots,\psi_{b_m}^m)))$$ iff $\M,w\Vdash\eta(\phi)$.
\end{proof}

%Ignored for First submission
%\nnote{New Lemma:}

\begin{lemma}\label{if homomorphic and back condition on atomic then also on composed}
    Let $\M$ and $\M'$ be $\tau$-models. Let $f:\M\to\M'$ be a mapping such that it satisfies the homomorphic condition (Definition~\ref{general homomorphism definition}) and the back condition (Definition~\ref{Bounded morphism - general}) w.r.t.\ the relations corresponding to basic modal terms. Then $f$ satisfies these conditions w.r.t.\ the relations corresponding to composed modal terms.
\end{lemma}

\begin{proof}
    Let $\M=(W,\{ R_{\alpha} \}_{\alpha \in \mtt},V)$ and $\M'=(W',\{ R'_{\alpha} \}_{\alpha \in \mtt},V')$ be $\tau$-models. Let $f:\M\to\M'$ be such that for all atomic $\alpha\in\mtt$, letting $\rho(\alpha) =m$, if $R_\alpha w_0w_1\ldots w_m$ then $R'_\alpha f(w_0)f(w_1)\ldots f(w_m)$,
    i.e., the homomorphic condition. Also let $f$ be such that all atomic $\alpha\in\mtt$ if $R'_\alpha f(w_0)w'_1\ldots w'_n$ then there exist $w_1,\ldots, w_n$ such that $R_\alpha ww_1\ldots w_n$ and $f(w_i) = w'_i$ (for $1\leq i \leq n)$, i.e., the back condition.
    \par Now, we will prove our original statement by induction on $\mtt$ --- Definition~\ref{cep of modal term def}. We already have our base case as an assumption. Now let $\alpha\in\mtt$ be a composed modality. I.e., letting $\rho(\alpha) = n$, $\alpha = \pi(\beta_1,\ldots,\beta_{\rho(\pi)})$ and assume that $R_\alpha w_0w_1\ldots w_n$ and that the homomorphic condition is satisfied for any modality with $cep$ value less than $cep(\alpha)$. Since $R_\alpha w_0w_1\ldots w_n$, we can rewrite this (letting $\rho(\pi) = m$ and $\rho(\beta_i) = b_i$) as $$R_{\pi(\beta_1,\ldots,\beta_m)} w_0w^1_1\ldots w^1_{b_1} \ldots w^m_1\ldots w^m_{b_m}.$$ Therefore there are some states $v_1,\ldots,v_m$ such that $R_\pi w_0v_1\ldots v_m$ \text{and} $R_{\beta_i} v_iw^i_1\ldots w^i_{b_i}$ for all $1\leq i \leq m$. Now since $cep(\pi)<cep(\alpha)$ and $cep(\beta_i)<cep(\alpha)$ for all $1\leq i \leq m$, then by our induction hypothesis $R'_\pi f(w_0)f(v_1)\ldots f(v_m)$ \text{ and } $R'_{\beta_i} f(v_i)f(w^i_1)\ldots f(w^i_{b_i})$ for all $1\leq i \leq m$. Hence $$R'_{\pi(\beta_1,\ldots,\beta_m)} f(w_0)f(w^1_1)\ldots f(w^1_{b_1}) \ldots f(w^m_1)\ldots f(w^m_{b_m}).$$ Therefore $R'_\alpha f(w_0)f(w_1)\ldots f(w_n)$, satisfying the homomorphic condition.
    \par Now, assume that the back condition is satisfied for any modality with $cep$ value less than $cep(\alpha)$. Also, assume $R'_\alpha f(w_0)w'_1\ldots w'_n$. This we can reindex the $w$'s (letting $\rho(\pi) = m$ and $\rho(\beta_i) = b_i$) as $$R'_{\pi(\beta_1,\ldots,\beta_m)} f(w_0)w'^1_1\ldots w'^1_{b_1} \ldots w'^m_1\ldots w'^m_{b_m}.$$ 
    Therefore there are some states $v'_1,\ldots,v'_m$ such that $R'_\pi f(w_0)v'_1\ldots v'_m$ \text{and} $R'_{\beta_i} v'_iw'^i_1\ldots w'^i_{b_i}$ for all $1\leq i\leq m$. We can see that $cep(\pi)<cep(\alpha)$ and so by our induction hypothesis we can say that there are states $v_1,\ldots,v_m\in W$ such that $$R_\pi w_0v_1\ldots v_m \text{ and } f(v_i) = v'_i$$ for all $1\leq i\leq m$. This means that $R'_{\beta_i} f(v_i)w'^i_1\ldots w'^i_{b_i}$ for all $1\leq i\leq m$. Since $cep(\beta_i)<cep(\alpha)$ for all $i$'s, we can apply the induction hypothesis on each $i$. Therefore, there are states $w^i_1,\ldots, w^i_{b_i}\in W$ such that $$R_{\beta_i} v_iw^i_1\ldots w^i_{b_i} \text{ and } f(w^i_j) = w'^i_j$$ for all $1\leq i\leq m$, $1\leq j\leq b_i$. Now, from $R_\pi w_0v_1\ldots v_m$ and $R_{\beta_i} v_iw^i_1\ldots w^i_{b_i}$ for each $1\leq i\leq m$, we can conclude that $$R_{\pi(\beta_1,\ldots,\beta_m)} w_0w^1_1\ldots w^1_{b_1} \ldots w^m_1\ldots w^m_{b_m}.$$ This shows that the back condition is satisfied since (after reindexing back) we have $R_\alpha w_0w_1\ldots w_n$ and we have found $w_1,\ldots,w_n$ such that $f(w_i) = w'_i$ for all $1\leq i\leq n$.
\end{proof}


\par This is all we need for the pure modal language, and we will now introduce a tableau system for this language.
